% GENERATED FILE. Edit canonical lessons, then run npm run generate:books.
% canonical-source-hash: fnv1a64:4f91b71c58622cc0
% canonical-lessons: GE-C481-allergie, GE-C481-allergisch, GE-C481-ausschlag, GE-C481-kratzen, GE-C481-eincremen, GE-R481-review-474-481

\chapter{Allergien und Ausschlag}
\label{ch:ge-allergien-und-ausschlag}
\hlchaptermodality{481}

\begin{chapteropening}
\textbf{By the end of this chapter:} \emph{I can talk about allergies and skin problems.}

Complete the last lesson of chapter 481: I can talk about allergies and skin problems.
\end{chapteropening}

\section[die Allergie]{die Allergie — allergy}

\label{lesson:GE-C481-allergie}



\begin{quote}
Before the new one: say the German for dizzy, then the German for circulation.
\end{quote}

\subsection*{You'll want to know: die Allergie}
\textbf{die Allergie} — \textquotedblleft{}allergy\textquotedblright{}.

\begin{quote}
Meine Tochter hat eine Allergie gegen Nüsse.
\end{quote}

\begin{quote}
\emph{My daughter has an allergy to nuts.}
\end{quote}

Plural: \emph{die Allergien}. \emph{Allergie gegen} + accusative.

\begin{cousinweb}
\emph{Allergie} is a modern learned word built from Greek \emph{állos} \textquotedblleft{}other\textquotedblright{} and \emph{érgon} \textquotedblleft{}work\textquotedblright{}: the body reacts in another way. EtymDB lists the German word without a source.
\end{cousinweb}

\begin{grammarlens}[title={What you've built}]
One more word for telling, explaining and describing at B1.
\end{grammarlens}

\subsection*{Your turn}
\begin{itemize}
\raggedright
  \item \emph{Say it:} \emph{die Allergie}
  \item \emph{Say it:} \emph{die Allergie}, once more
  \item \emph{Say it:} a sentence of your own with \emph{die Allergie}
  \item \emph{Recall:} say the German for dizzy, then the German for circulation, then say \emph{die Allergie} again
\end{itemize}

\begin{tcolorbox}[breakable,colback=teal!4,colframe=teal!35!black,title={Before you move on}]
What does die Allergie mean? (\textquotedblleft{}Allergy\textquotedblright{}.) Say it once more.
\end{tcolorbox}
\section[allergisch]{allergisch — allergic}

\label{lesson:GE-C481-allergisch}



\begin{quote}
Before the new one: say the German for circulation, then the German for allergy.
\end{quote}

\subsection*{You'll want to know: allergisch}
\textbf{allergisch} — \textquotedblleft{}allergic\textquotedblright{}.

\begin{quote}
Ich bin allergisch gegen Katzen.
\end{quote}

\begin{quote}
\emph{I'm allergic to cats.}
\end{quote}

\emph{Allergisch gegen} + accusative, also as a joke: \emph{allergisch gegen Montage}.

\begin{cousinweb}
\emph{allergisch} is built on \emph{Allergie} with the ending \emph{-isch} (EtymDB); \emph{Allergie} is a modern learned word from Greek \emph{állos} \textquotedblleft{}other\textquotedblright{} and \emph{érgon} \textquotedblleft{}work\textquotedblright{}.
\end{cousinweb}

\begin{grammarlens}[title={What you've built}]
One more word for telling, explaining and describing at B1.
\end{grammarlens}

\subsection*{Your turn}
\begin{itemize}
\raggedright
  \item \emph{Say it:} \emph{allergisch}
  \item \emph{Say it:} \emph{allergisch}, once more
  \item \emph{Say it:} a sentence of your own with \emph{allergisch}
  \item \emph{Recall:} say the German for circulation, then the German for allergy, then say \emph{allergisch} again
\end{itemize}

\begin{tcolorbox}[breakable,colback=teal!4,colframe=teal!35!black,title={Before you move on}]
What does allergisch mean? (\textquotedblleft{}Allergic\textquotedblright{}.) Say it once more.
\end{tcolorbox}
\section[der Ausschlag]{der Ausschlag — rash}

\label{lesson:GE-C481-ausschlag}



\begin{quote}
Before the new one: say the German for allergy, then the German for allergic.
\end{quote}

\subsection*{You'll want to know: der Ausschlag}
\textbf{der Ausschlag} — \textquotedblleft{}rash\textquotedblright{}.

\begin{quote}
Nach der neuen Creme bekam sie einen Ausschlag.
\end{quote}

\begin{quote}
\emph{After the new cream she got a rash.}
\end{quote}

Plural: \emph{die Ausschläge}. \emph{Den Ausschlag geben} — to tip the balance.

\begin{cousinweb}
\emph{Ausschlag} is built on \emph{aus} \textquotedblleft{}out\textquotedblright{} and \emph{Schlag} \textquotedblleft{}blow\textquotedblright{}, inherited from Old High German \emph{slag}, from Proto-Germanic \emph{slagiz} (EtymDB): something that breaks out on the skin.
\end{cousinweb}

\begin{grammarlens}[title={What you've built}]
One more word for telling, explaining and describing at B1.
\end{grammarlens}

\subsection*{Your turn}
\begin{itemize}
\raggedright
  \item \emph{Say it:} \emph{der Ausschlag}
  \item \emph{Say it:} \emph{der Ausschlag}, once more
  \item \emph{Say it:} a sentence of your own with \emph{der Ausschlag}
  \item \emph{Recall:} say the German for allergy, then the German for allergic, then say \emph{der Ausschlag} again
\end{itemize}

\begin{tcolorbox}[breakable,colback=teal!4,colframe=teal!35!black,title={Before you move on}]
What does der Ausschlag mean? (\textquotedblleft{}Rash\textquotedblright{}.) Say it once more.
\end{tcolorbox}
\section[kratzen]{kratzen — to scratch}

\label{lesson:GE-C481-kratzen}



\begin{quote}
Before the new one: say the German for allergic, then the German for rash.
\end{quote}

\subsection*{You'll want to know: kratzen}
\textbf{kratzen} — \textquotedblleft{}to scratch\textquotedblright{}.

\begin{quote}
Die Katze hat mich am Arm gekratzt.
\end{quote}

\begin{quote}
\emph{The cat scratched my arm.}
\end{quote}

\emph{Sich kratzen} — to scratch yourself: \emph{Er kratzt sich am Kopf.}

\begin{cousinweb}
\emph{kratzen} is inherited from Middle High German \emph{kratzen} and Old High German \emph{krazzōn}, from Proto-Germanic \emph{krattōną} \textquotedblleft{}to scratch\textquotedblright{} (EtymDB).
\end{cousinweb}

\begin{grammarlens}[title={What you've built}]
One more word for telling, explaining and describing at B1.
\end{grammarlens}

\subsection*{Your turn}
\begin{itemize}
\raggedright
  \item \emph{Say it:} \emph{kratzen}
  \item \emph{Say it:} \emph{kratzen}, once more
  \item \emph{Say it:} a sentence of your own with \emph{kratzen}
  \item \emph{Recall:} say the German for allergic, then the German for rash, then say \emph{kratzen} again
\end{itemize}

\begin{tcolorbox}[breakable,colback=teal!4,colframe=teal!35!black,title={Before you move on}]
What does kratzen mean? (\textquotedblleft{}To scratch\textquotedblright{}.) Say it once more.
\end{tcolorbox}
\section[eincremen]{eincremen — to put cream on}

\label{lesson:GE-C481-eincremen}



\begin{quote}
Before the new one: say the German for rash, then the German for to scratch.
\end{quote}

\subsection*{You'll want to know: eincremen}
\textbf{eincremen} — \textquotedblleft{}to put cream on\textquotedblright{}.

\begin{quote}
Bevor wir an den Strand gehen, cremen wir die Kinder gut ein.
\end{quote}

\begin{quote}
\emph{Before the beach we put plenty of cream on the children.}
\end{quote}

Separable: \emph{ich creme ein}. Also reflexive: \emph{sich eincremen}.

\begin{cousinweb}
\emph{eincremen} is built on \emph{ein} \textquotedblleft{}in\textquotedblright{} and \emph{Creme}, which was borrowed from French \emph{crème} \textquotedblleft{}cream\textquotedblright{} (EtymDB); EtymDB has no entry for the verb.
\end{cousinweb}

\begin{grammarlens}[title={What you've built}]
One more word for telling, explaining and describing at B1.
\end{grammarlens}

\subsection*{Your turn}
\begin{itemize}
\raggedright
  \item \emph{Say it:} \emph{eincremen}
  \item \emph{Say it:} \emph{eincremen}, once more
  \item \emph{Say it:} a sentence of your own with \emph{eincremen}
  \item \emph{Recall:} say the German for rash, then the German for to scratch, then say \emph{eincremen} again
\end{itemize}

\begin{tcolorbox}[breakable,colback=teal!4,colframe=teal!35!black,title={Before you move on}]
What does eincremen mean? (\textquotedblleft{}To put cream on\textquotedblright{}.) Say it once more.
\end{tcolorbox}
\section[(dialogue)]{Words from chapters 474–481 — a review}

\label{lesson:GE-R481-review-474-481}



\begin{quote}
Say the German for to be amazed and astonished.
\end{quote}

\subsection*{Your turn}
Say these aloud:

\begin{itemize}
\raggedright
  \item to be amazed, astonished, confused, curiosity, speechless
  \item to boast, to show off, show-off, modest, vain
  \item stubborn, persistent, moody, irritable, sensitive
  \item brave, cowardly, coward, reckless, to dare
  \item remorse, to regret, to forgive, to insult, to hurt
  \item to catch a cold, head cold, cough syrup, sweat, goose bumps
  \item nausea, to be sick, diarrhoea, dizzy, circulation
  \item allergy, allergic, rash, to scratch, to put cream on
\end{itemize}

\begin{tcolorbox}[breakable,colback=teal!4,colframe=teal!35!black,title={Before you move on}]
To be amazed? (\textbf{staunen}.) Astonished? (\textbf{erstaunt}.) Confused? (\textbf{verwirrt}.) Curiosity? (\textbf{die Neugier}.) Speechless? (\textbf{sprachlos}.) To boast? (\textbf{prahlen}.) To show off? (\textbf{angeben}.) Show-off? (\textbf{der Angeber}.) Modest? (\textbf{bescheiden}.) Vain? (\textbf{eitel}.) Stubborn? (\textbf{stur}.) Persistent? (\textbf{hartnäckig}.) Moody? (\textbf{launisch}.) Irritable? (\textbf{gereizt}.) Sensitive? (\textbf{empfindlich}.) Brave? (\textbf{tapfer}.) Cowardly? (\textbf{feige}.) Coward? (\textbf{der Feigling}.) Reckless? (\textbf{leichtsinnig}.) To dare? (\textbf{trauen}.) Remorse? (\textbf{die Reue}.) To regret? (\textbf{bereuen}.) To forgive? (\textbf{verzeihen}.) To insult? (\textbf{beleidigen}.) To hurt? (\textbf{kränken}.) To catch a cold? (\textbf{erkälten}.) Head cold? (\textbf{der Schnupfen}.) Cough syrup? (\textbf{der Hustensaft}.) Sweat? (\textbf{der Schweiß}.) Goose bumps? (\textbf{die Gänsehaut}.) Nausea? (\textbf{die Übelkeit}.) To be sick? (\textbf{übergeben}.) Diarrhoea? (\textbf{der Durchfall}.) Dizzy? (\textbf{schwindlig}.) Circulation? (\textbf{der Kreislauf}.) Allergy? (\textbf{die Allergie}.) Allergic? (\textbf{allergisch}.) Rash? (\textbf{der Ausschlag}.) To scratch? (\textbf{kratzen}.) To put cream on? (\textbf{eincremen}.)
\end{tcolorbox}
